\section{Modelo de Inferência}

Para o armazenamento e processamento das ontologias neste trabalho, foram estudadas algumas linguagens de representação de conhecimento.

% CycL
A linguagem \textbf{CycL} é a linguagem utilizada pelo o projeto de \emph{Cyc}\cite{lenat89}. O Cyc é um projeto de inteligência artificial que tem o objetivo de montar uma ontologia e base de conhecimento a partir do senso comum, com o propósito de permitir que aplicações de inteligência artificial realizem raciocínio semelhante ao humano. A linguagem CycL é uma linguagem declarativa baseada em lógica de primeira ordem clássica, com extensões para operadores modais e lógica de ordem superior. A base de conhecimento é dividida em micro-teorias (Mt), que são coleções de conceitos e fatos tipicamente pertencentes a um domínio particular de conhecimentos e, ao contrário da base de conhecimento como um todo, cada micro-teoria deve estar livre de contradições.

% RDF/S
\textbf{RDF} (\emph{Resource Description Framework}) é uma linguagem utilizada para representação de metadados a respeito de recursos na Web\cite{rdf}. No entanto, o conceito de ``recurso da Web'' pode ser generalizado, de forma que o RDF pode ser utilizado para representar informações a respeito de coisas que podem ser identificadas na Web, mesmo que elas não possam ser carregadas da Web. Recursos podem ser recursos eletrônicos, como arquivos, ou conceitos, como pessoas. O RDF é complementado pelo \textbf{RDF/S} (\emph{RDF/Schema}), uma extensão semântica do RDF. O RDF/S é uma linguagem de descrição de vocabulário que provê mecanismos para descrever grupos de recursos relacionados e os relacionamentos entre estes recursos\cite{rdfs}.

% DAML+OIL
O \textbf{DAML+OIL} é uma linguagem de representação de conhecimento baseada nas linguagens DAML e OIL. O DAML (\emph{DARPA Agent Markup Language}) foi desenvolvido pelo Centro de Pesquisas Avançadas da Defesa Americana e é uma linguagem de marcação de agentes baseada no RDF. Já o OIL (\emph{Ontology Inference Layer}) é uma linguagem com os mesmos objetivos, criada pelo por um grupo de universidades europeias. Ambos os esforços foram unidos na criação do DAML+OIL, que por sua vez deu origem ao OWL\cite{daconta03}.

% OWL
O \textbf{OWL} é a linguagem de ontologia mais expressiva definida ou em definição para a Web Semântica\cite{daconta03}. O OWL utiliza URI\footnote{\emph{Uniform Resource Identifier}} para os nomes, e a estrutura de descrição disponibilizada pelo RDF. O OWL está uma camada acima do RDF e do RDF/S, e adiciona mais vocabulário para descrever propriedades e classes: entre outros, relações entre classes, cardinalidade, igualdade, tipagem mais avançada de propriedades, características de propriedades e classes enumeradas\cite{owl-guide}. O OWL está disponível em três sublinguages: \emph{OWL Lite}, \emph{OWL DL} e \emph{OWL Full}, cada uma aumentando o nível de complexidade e expressividade. Um grande número de ferramentas com capacidade para edição e raciocínio semântico com suporte para OWL está disponível no mercado.

% Escolhido - OWL
Devido à expressividade e ao suporte de ferramentas externas, o OWL será a linguagem de ontologias utilizado pelo eProfile para modelagem da ontologia do perfil.

% Raciocinadores
Existem diversos raciocinadores que oferecem suporte a OWL, entre os quais estão:

\begin{description}

\item[Pellet:] é um racionador com suporte a OWL DL, que oferece bom desempenho e um \emph{middleware} extensível. O Pellet é escrito em Java pode ser utilizado para raciocinar com tipos de dados definidos pelo usuário e provê suporte para depuração em ontologias\cite{sirin05}.

\item[FaCT++:] é um raciocinador de lógica proposicional criado como uma plataforma para experimentar com novos algoritmos de \emph{tableaux} semântico e técnicas de otimização. Ele incorpora a maior parte das técnicas de otimização padrão, mas inclui algumas novas\cite{tsarkov06}.

\item[HermiT:] é o primeiro raciocinador OWL que se baseia em um cálculo \emph{hypertableaux}, que provê desempenho superior a qualquer algoritmo anterior\cite{motik09}. O HermiT é disponibilizado como um plugin do Protegé, como uma ferramenta de linha de comando ou como uma biblioteca Java.

\end{description}
